% ==============================================================================
% Plantilla Institucional de Trabajo Terminal — UPIIZ IPN
% Archivo: capitulos/04-marco-teorico.tex
% ==============================================================================
% LÍMITE NORMATIVO INSTITUCIONAL: Máximo 3 páginas.
% ==============================================================================

\chapter{Marco Teórico}
\label{cap:marco-teorico}

En este capítulo se presentan los fundamentos físico-matemáticos indispensables que sustentan el modelado dinámico, el análisis de estabilidad, las leyes de control y las técnicas de procesamiento empleadas en el proyecto \cite{ogata2010modern}.

% ------------------------------------------------------------------------------
\section{Modelado físico-matemático del sistema}
\label{sec:marco-modelado}

Describa las ecuaciones diferenciales que gobiernan el comportamiento del sistema mecatrónico. Por ejemplo, para un sistema electromecánico lineal invariante en el tiempo continuo, la representación general en espacio de estados se formula como:
\begin{align}
    \dot{\mathbf{x}}(t) &= \mathbf{A}\mathbf{x}(t) + \mathbf{B}\mathbf{u}(t) \label{eq:espacio-estados-continuo} \\
    \mathbf{y}(t) &= \mathbf{C}\mathbf{x}(t) + \mathbf{D}\mathbf{u}(t) \label{eq:salida-continuo}
\end{align}
donde $\mathbf{x}(t) \in \mathbb{R}^n$ es el vector de variables de estado, $\mathbf{u}(t) \in \mathbb{R}^m$ es el vector de entradas de control, $\mathbf{y}(t) \in \mathbb{R}^p$ es el vector de salidas medidas, y $\mathbf{A}, \mathbf{B}, \mathbf{C}, \mathbf{D}$ son las matrices del sistema.

% ------------------------------------------------------------------------------
\section{Análisis de controlabilidad y leyes de control}
\label{sec:marco-control}

La controlabilidad del sistema garantiza que es posible transferir el estado del sistema desde cualquier condición inicial hasta el origen en tiempo finito mediante una ley de control admisible. La matriz de controlabilidad $\mathcal{C}$ se define como:
\begin{equation}
    \mathcal{C} = \begin{bmatrix} \mathbf{B} & \mathbf{A}\mathbf{B} & \mathbf{A}^2\mathbf{B} & \cdots & \mathbf{A}^{n-1}\mathbf{B} \end{bmatrix}, \quad \text{rank}(\mathcal{C}) = n
    \label{eq:matriz-controlabilidad-marco}
\end{equation}

La ley de control por realimentación de estados con precompensación para seguimiento de referencias $\mathbf{r}(t)$ se expresa mediante:
\begin{equation}
    \mathbf{u}(t) = -\mathbf{K}\mathbf{x}(t) + \mathbf{N}\mathbf{r}(t)
    \label{eq:ley-control-marco}
\end{equation}
donde $\mathbf{K}$ representa la matriz de ganancias calculada para satisfacer las especificaciones dinámicas de tiempo de establecimiento y sobrepaso porcentual.

Recuerde expresar todas las magnitudes físicas utilizando el paquete \texttt{siunitx}, por ejemplo: tensión de \qty{12.0}{\volt}, corriente de \qty{2.5}{\ampere}, frecuencia de muestreo de \qty{1.0}{\kilo\hertz} y par motor de \qty{0.45}{\newton\meter}.
